\begin{proposition}[Far-bridge incidence majorants]
\label{prop:far-bridge-incidence}
Fix an odd prime~$p$.  Let $\pi_p$ be the projective Ap\'ery orbit on
$\{0,\ldots,p-1\}$, and let $\Gamma_p(q)$ be a reflection-invariant
collection of filtered four-consecutive-occurrence windows as in
Proposition~\ref{prop:projective-variance-reduction}.  Assume that every
selected window has span at most
\[
 H_p=\lfloor\sqrt p\rfloor.
\]
For each ordered separated pair of distinct reflection orbits, make the
deterministic choice of oriented representatives used in
Proposition~\ref{prop:near-bridge-energy}, and write its bridge as~$G$.
Let $E_p^{\mathrm{far}}(K)$ count the pairs for which $G>K$, where
$1\le K<p$ is an integer.

Define the conditioned incidence
\[
 B_p(K)=
 \sum_{\substack{q\in\mathbf P^1(\mathbf F_p)\\
                   \gamma=(x_0,x_1,x_2,x_3)\in\Gamma_p(q)}}
 \#\{z:x_3+K<z<p,\ \pi_p(z)=\pi_p(x_3)\},
\]
and the nonwrapping triple incidence
\[
 A_p^{\mathrm{nw}}(K)=
 \#\left\{(x,y,z):
 \begin{array}{c}
 0\le x<y<z<p,\quad 2\le y-x\le H_p,\quad z-y>K,\\
 \pi_p(x)=\pi_p(y)=\pi_p(z)
 \end{array}
 \right\}.
\]
Then
\begin{equation}
 E_p^{\mathrm{far}}(K)\le B_p(K)\le A_p^{\mathrm{nw}}(K).
\label{eq:far-incidence-hierarchy}
\end{equation}
Moreover, if $N_h$ denotes the positional-gap continuant, then
\begin{equation}
\begin{split}
 A_p^{\mathrm{nw}}(K)
 ={}&\sum_{s=2}^{H_p}\ \sum_{G=K+1}^{p-1-s}
 \#\bigl\{x\in\mathbf Z:
 0\le x\le p-1-s-G,\\[-2pt]
 &\hspace{105pt}
 N_s(x)\equiv N_G(x+s)\equiv0\pmod p\bigr\}.
\end{split}
\label{eq:far-incidence-continuants}
\end{equation}
In particular, it is bounded above by the unrestricted affine-root mass
\begin{equation}
 \widetilde A_p(K)=
 \sum_{s=2}^{H_p}\ \sum_{G=K+1}^{p-1-s}
 \#\{u\in\mathbf F_p:
       N_s(u)=N_G(u+s)=0\}.
\label{eq:far-incidence-affine}
\end{equation}
The summands in~\eqref{eq:far-incidence-affine} count actual
$\mathbf F_p$-roots; they are not degrees of polynomial greatest common
divisors.
\end{proposition}

\begin{proof}
Take an ordered separated pair counted by
$E_p^{\mathrm{far}}(K)$.  Its deterministic representatives have the form
\[
 \gamma=(x_0,x_1,x_2,x_3),\qquad
 \delta=(y_0,y_1,y_2,y_3),
 \qquad x_3<y_0,\qquad y_0-x_3>K.
\]
Both windows belong to the same projective fiber, so
$\pi_p(y_0)=\pi_p(x_3)$.  Map the ordered quotient pair to
$(\gamma,y_0)$.

This map is injective.  Indeed, $\gamma$ determines the first reflection
orbit.  Given $y_0$, the next three occurrences of
$\pi_p(y_0)$ are uniquely determined by the ordered fiber.  Hence there is
at most one four-consecutive-occurrence window beginning at~$y_0$, and, when
the image comes from a selected pair, it is precisely~$\delta$.  Thus the
second reflection orbit is recovered as well.  The image is among the
incidences counted by~$B_p(K)$, proving the first inequality in
\eqref{eq:far-incidence-hierarchy}.  No factor of two is lost here: the
canonical choice assigns one oriented pair to each ordered pair of quotient
orbits, whereas $B_p(K)$ already sums over all selected orientations.

For an incidence $(\gamma,z)$ counted by~$B_p(K)$, set
\[
 x=x_0,\qquad y=x_3.
\]
The selected-window hypotheses give
$2\le y-x\le H_p$, and
\[
 \pi_p(x)=\pi_p(y)=\pi_p(z),\qquad z-y>K.
\]
The map $(\gamma,z)\mapsto(x,y,z)$ is injective, because the two internal
indices of a four-consecutive-occurrence window are the first two
occurrences of $\pi_p(x)$ strictly between its endpoints.  This proves
$B_p(K)\le A_p^{\mathrm{nw}}(K)$.

It remains only to reparametrize the latter set.  Put
\[
 s=y-x,\qquad G=z-y.
\]
Then
\[
 2\le s\le H_p,\qquad K+1\le G\le p-1-s,\qquad
 0\le x\le p-1-s-G.
\]
All two-time intervals are nonwrapping.  The exact gap criterion therefore
gives
\[
 \pi_p(x)=\pi_p(x+s)
 \iff N_s(x)\equiv0\pmod p
\]
and
\[
 \pi_p(x+s)=\pi_p(x+s+G)
 \iff N_G(x+s)\equiv0\pmod p.
\]
This proves~\eqref{eq:far-incidence-continuants}.  Dropping the interval
condition on the standard representative of~$x$ gives
$A_p^{\mathrm{nw}}(K)\le\widetilde A_p(K)$.

Finally, a degree such as
$\deg\gcd_{\mathbf F_p[x]}(N_s(x),N_G(x+s))$ also counts common irreducible
factors having no root in~$\mathbf F_p$.  It is therefore only an upper
bound for the corresponding summand in
\eqref{eq:far-incidence-affine}, not the exact observable in this
proposition.
\end{proof}

\begin{remark}[Scale of the remaining arithmetic input]
By Lemma~\ref{lem:nonvanish}, for $s,G<p$ neither gap polynomial vanishes
identically modulo~$p$.  The actual common-root count is at most the degree
of the greatest common divisor, and hence at most $3(s-1)$.  Consequently
the elementary degree estimate gives only
\[
 \widetilde A_p(K)
 \le\sum_{s\le H_p}\sum_{G<p}3(s-1)\ll pH_p^2\ll p^2.
\]
After summing over $X<p\le2X$, this is
$O(X^3/\log X)$, one full power of~$X$ above the required separated-energy
budget.  Thus~\eqref{eq:far-incidence-hierarchy} is useful only together
with an average arithmetic estimate for the actual affine roots, or with
the sharper conditioning retained in~$B_p(K)$.
\end{remark}
